\chapter{CONCLUSÃO}

% Parágrafo introdutório: resumir brevemente o que o capítulo apresenta
% (síntese do trabalho, respostas às hipóteses, considerações práticas, limitações e trabalhos futuros)

\section{Síntese do Trabalho}

% Resumir em 2–3 parágrafos:
% - o objetivo do laboratório e a motivação
% - os modelos comparados (RF, DNN, Autoencoder) e os datasets (CICIDS2017, LycoS-IDS2017)
% - as linhas de melhoria investigadas (A–G)

\section{Respostas às Hipóteses}

% Retomar as hipóteses do Capítulo 1 e responder explicitamente se foram confirmadas ou refutadas.
%
% Hipótese 1 (DL superior em ataques complexos e classes de baixo suporte):
% — REFUTADA no contexto estudado: o RF superou a DNN em todos os cenários,
%   incluindo classes raras. A DNN melhorou com o LycoS-IDS2017, mas nunca
%   igualou o RF em macro F1.
%
% Hipótese 2 (RF mais eficiente computacionalmente):
% — CONFIRMADA: o RF foi 5 a 6 vezes mais rápido na inferência em todos os cenários.

\section{Considerações Práticas}

% Discutir qual abordagem é mais adequada para cada cenário de uso real.
% Sugestões de tópicos (ver TODO 4 no arquivo de melhorias):
% - estabilidade do RF vs. custo e fragilidade da DNN
% - impacto da qualidade do dataset (LycoS vs. CICIDS2017)
% - quando usar o Autoencoder (ataques zero-day, sem rótulos)
% - trade-off entre granularidade (supervisionado) e generalização (não supervisionado)
% - viabilidade de cada abordagem em ambientes com recursos limitados

\section{Limitações do Trabalho}

Este trabalho apresenta algumas delimitações de escopo que devem ser consideradas na interpretação dos resultados.

Os experimentos utilizam datasets públicos gerados em ambiente controlado, o que pode não refletir toda a diversidade de tráfego encontrada em redes corporativas reais \cite{cantone2024}. A avaliação foi realizada de forma offline sobre conjuntos de dados fixos, sem testes em tempo real ou em fluxo contínuo. Os modelos supervisionados dependem de exemplos rotulados de cada classe de ataque para o treinamento, enquanto o Autoencoder, por ser treinado apenas com tráfego benigno, reduz essa dependência ao custo de não distinguir entre tipos de ataque. Por fim, os datasets utilizados não segregam fluxos por protocolo de criptografia, de modo que a eficácia dos modelos em redes com alto volume de tráfego TLS/HTTPS não foi avaliada neste trabalho.

\section{Trabalhos Futuros}

Os resultados e as delimitações de escopo deste trabalho apontam para diversas direções que podem ampliar e aprofundar os experimentos realizados.

A interpretabilidade dos modelos constitui uma extensão imediata. Os modelos de DNN e Autoencoder investigados operam como caixas-pretas, sem explicar quais atributos motivaram cada decisão. A integração de mecanismos de \textit{Explainable AI} (XAI), como SHAP (\textit{SHapley Additive exPlanations}) ou LIME, permitiria identificar as \textit{features} mais discriminativas para cada tipo de ataque e auxiliar os analistas na interpretação dos alertas gerados \cite{ravindran2025}.

A avaliação em tempo real representa outro caminho natural. Os experimentos foram realizados em modo \textit{batch} sobre datasets fixos; a extensão para fluxo contínuo introduziria requisitos de latência e adaptação a mudanças na distribuição do tráfego. A arquitetura de dois estágios baseada em Filtro de Bloom proposta por Rego e Nunes \cite{rego_nunes} é um exemplo de abordagem que concilia baixa latência com o poder discriminativo de redes neurais. Nessa mesma direção, a pesquisa de modelos \textit{lightweight} adequados para implantação em dispositivos de borda (\textit{edge computing}), com restrições severas de memória e processamento, representa um problema prático relevante.

A generalização entre datasets é uma das questões mais abertas na área. Cantone et al. \cite{cantone2024} demonstraram que modelos treinados em um dataset tendem a fracassar ao serem avaliados em outros. A avaliação \textit{cross-dataset}, com modelos treinados em \textit{CICIDS2017} ou \textit{LycoS-IDS2017} e avaliados em conjuntos distintos como o \textit{CIC-IDS-2018} ou o \textit{UNSW-NB15}, permitiria medir o grau de generalização real dos modelos desenvolvidos.

Cenários com tráfego criptografado constituem uma direção independente. Os datasets utilizados não segmentam os fluxos por protocolo de criptografia, e a análise de redes com alto volume de tráfego TLS/HTTPS, onde os atributos de \textit{payload} são inacessíveis, exigiria \textit{features} derivadas exclusivamente de metadados e estatísticas de fluxo.

Por fim, a aprendizagem federada (\textit{federated learning}) representa uma alternativa promissora para cenários onde a centralização de dados de tráfego de rede é inviável por razões de privacidade ou regulatórias. Nesse paradigma, cada organização treina um modelo local com seus próprios dados e compartilha apenas as atualizações de pesos, preservando a confidencialidade das informações de rede.
